\appendix
\section{Auditable normal forms for theorem statements}

\subsection{Deterministic finite construction}

A deterministic theorem should be presented as
\[
  \forall x\quad \Good(x)=1\longrightarrow
  \Post(x,F(x))=1,
\]
where \(F\), \(\Good\), and \(\Post\) are fixed primitive-recursive symbols or finite-observation predicates.  The artifact names the trace, invariant, fuel, and local derivations.

\subsection{Precision-indexed construction}

A represented theorem should be presented as
\[
  \forall p\quad H(p)\longrightarrow
  \forall k\quad \Obs(p,k,A(p,k))=1.
\]
The proof keeps \(p,k\) free, computes a finite query plan, proves the observation, and only then universally introduces \(k\).  The interface lists every modulus, magnitude bound, and guard consumed.

\subsection{Relational finite-use theorem}

A multivalued producer should expose
\[
  \forall x,h\quad \mathsf{Evidence}(x,h)=1\longrightarrow
  \Viable(x,h,\mathsf{start}(x))=1.
\]
A consumer theorem concludes a successful finite transcript, a quotient-invariant observation, or a diameter-stable output.  It must not silently replace finite viability by one infinite coherent selector.

\subsection{Alternating bounded contract}

A finite challenge theorem should expose both sides:
\[
 \CertCheck_A(x,1,c)=1\to A(x),\qquad
 \CertCheck_A(x,0,c)=1\to\neg A(x).
\]
The canonical evaluator produces one of the two certificates.  This makes failed finite stages auditable and separates logical failure from a missing external producer.

\subsection{Conditional certificate theorem}

For an external certificate \(c\), state
\[
  \forall x,c\quad \CertCheck(x,c)=1\longrightarrow
  \Phi(x,\mathsf{out}(c)).
\]
Record separately whether a producer is known for all admissible \(x\), for a special family, or only experimentally.

\subsection{Exact branch under a gap}

An exact selector should have the guarded form
\[
  \mathsf{Gap}(x,m)=1\longrightarrow
  \mathsf{ExactBranch}(x,\mathsf{select}(x,m))=1.
\]
Without a gap, use overlapping branches, a finite relation, a cluster, a projective class, or a stable observation.

\subsection{Height and complete-list theorem}

A Diophantine enumeration theorem should name the effective producer:
\[
  \mathsf{CurveHyp}(x)\longrightarrow
  \forall z\bigl(P(x,z)=1\to h(z)\leq B(x)\bigr),
\]
with concrete primitive-recursive \(B\), or else accept a supplied height certificate.  Only after this premise is available does the generic bounded enumerator yield a complete list.

\section{Notation and implementation conventions}

\begin{center}
\begin{tabularx}{\textwidth}{@{}lX@{}}
\toprule
Notation & Meaning \\
\midrule
\code{PR} & fixed primitive-recursive definition, never an unclocked universal evaluator \\
\(u,x,p\) & paired numerical input, parameter, or certificate codes \\
\(k\) & requested output precision \\
\(h,r\) & finite interaction horizon or remaining resource \\
\(s\) & encoded machine/consumer state \\
\(\Viable(x,r,s)\) & backward Boolean saying a successful finite continuation exists \\
\(\Local(x,r,s,a)\) & local admissibility of answer \(a\) \\
\(E\) & observational equivalence, not necessarily literal equality \\
\(\delta\) & output diameter or error bound \\
\(\Prf\ \PRAAx\ \Gamma\ \varphi\) & object-level derivability in ordinary iPRA \\
\bottomrule
\end{tabularx}
\end{center}

All finite sequences and graphs may use any standard primitive-recursive coding with explicit bounds.  Implementations should nonetheless share one canonical repository encoding to avoid repeated conversion theorems.  Rational inequalities are reduced to signed-integer cross products with positive denominators.  Every order decision is therefore an exact decision on finite codes, not an order decision on an arbitrary represented real.
